\scanpage{9}\mspage{6}

\[
 L^{n-i}\mid M^i : M^i\longrightarrow M^{2n-i}\quad\text{est un isom.}
 \tag*{$(6.1)_i$}
\]
Ceci implique déjà $(4.1)_i$ et $(4.2)_i$ donc $(3.1)_i$ et $(3.2.)_i$, \inserted{donc l'\uline{unicité} des $\pi_i$ en termes de la structure graduée de $\unsure{E|\mathcal E}^*$.} Supposons de plus

$(6.2)_i$ L'inverse de l'isom (6.1) \deleted{\unsure{dans $E$|$E$}\illegible{1}} \inserted{est} induit par un \unsure{élmt|élt} de $E^{-(2n-2i)}$.

\deleted{Ceci} L'ensemble des deux conditions (6.1) et (6.2). équivaut [en termes intrinsèques à $E^*$, muni de \deleted{\unsure{$E$}} $L\in E^2$, si on suppose déjà les $\pi_i\in E^*$]
\[
 \left\{\begin{aligned}
 &\exists\,\unsure{W_i|w_i}\in E^{-(2n-2i)}\quad\text{tel qu'on ait}\\
 &(w_i L^{n-i}-1)\pi_i=0,\qquad (L^{n-i}w_i-1)\pi_{2n-i}=0
 \end{aligned}\right.
 \tag*{$(6.3.)_i$}
\]
Cette condition à son tour (indépendamment de la donnée préliminaire d'un $M^*$) implique $(4.4)_i$, \deleted{et on en déduit $(3.1)_i$ et $(3.2)_i$} donc \inserted{en vertu des n\textsuperscript{os} 4, 5} \inserted{$\pi_i\in E$ et} l'\uline{unicité} des $\pi_i$ en termes des conditions \inserted{(1.1) et} (1.3).

Revenons à $M^*$, et $L$ satisfaisant à $(6.1)_i$. Nous avons vu comment \deleted{\unsure{la}} cette condition implique la \uline{décomposition de} $M^*$ en termes \deleted{\unsure{nou}} de morceaux primitifs
\[
 \left\{\begin{aligned}
 M^i&\simeq P^i+LP^{i-2}+L^2P^{i-4}\ \unsure{-|\cdots}\\
 M^{2n-i}&\simeq L^{n-i}P^i+L^{n-i+1}P^{i-2}\ \unsure{-|\cdots}
 \end{aligned}\ \right|\quad\text{si }i\le n
 \tag*{(6.4)}
\]
On désigne par $\pi_{i,\alpha}$ la projection de $M^i$ sur $L^\alpha P^{i-2\alpha}$, par $\pi_{2n-i,\alpha}$ celle de $M^{2n-i}$ sur $L^{\unsure{ni+\alpha|n-i+\alpha}}P^{i-2\alpha}$.
% archive sheet number: 8
